% !TEX root = ../main.tex
\begin{figure}[H]
  \centering
  \resizebox{\textwidth}{!}{%
    \begin{tikzpicture}[node distance=0.55cm and 1cm]
      % Fase 0
      \node[nodoDiag] (web) {Web pública\\EPSJ};
      \node[nodoDiag, below=of web] (grados) {\texttt{grados.json}};
      \node[nodoDiag, below=of grados] (chunks) {\texttt{chunks.json}};
      % Fase 1
      % Actualizado el 30/07/2026: el modelo de incrustaciones dejo de ser una
      % caja generica cuando el ADR-0003 lo fijo, asi que la figura lo nombra.
      % Actualizado el 13/08/2026: la base de datos deja de ser una caja generica
      % por el mismo motivo que el modelo, ahora que el ADR-0004 la fija.
      \node[nodoDiag, right=3.4cm of web] (emb) {Modelo de\\\textit{embeddings}\\{\scriptsize multilingual-e5-small}};
      \node[nodoDiag, cylinder, shape border rotate=90, aspect=0.15, below=of emb] (bd) {Base de datos\\vectorial\\{\scriptsize LanceDB}};
      % Fase 2
      \node[nodoDiag, right=3cm of emb] (preg) {Pregunta\\del usuario};
      \node[nodoDiag, below=of preg] (rec) {Recuperación\\de fragmentos};
      \node[nodoDiag, below=of rec] (llm) {LLM local};
      \node[nodoDiag, below=of llm] (resp) {Respuesta\\fundamentada};
      % Fase 3
      \node[nodoDiag, right=3.2cm of rec] (chat) {Interfaz de chat\\sin registro};
      % Nodos invisibles que "miden" cada título (mismo texto y fuente que la
      % etiqueta real) para que el fit incluya su ancho: un label de TikZ no
      % forma parte del cálculo del área del nodo, así que sin esto el título
      % puede sobresalir de la caja aunque esta se ensanche.
      \node[font=\small, inner sep=0pt, text opacity=0, above=1pt of web] (f0medida) {Fase 0 --- Colección de datos};
      \node[font=\small, inner sep=0pt, text opacity=0, above=1pt of emb] (f1medida) {Fase 1 --- Indexación};
      \node[font=\small, inner sep=0pt, text opacity=0, above=1pt of preg] (f2medida) {Fase 2 --- \textit{Pipeline} RAG};
      \node[font=\small, inner sep=0pt, text opacity=0, above=1pt of chat] (f3medida) {Fase 3 --- Aplicación web};
      % Cajas de fase (en la capa de fondo, para no tapar los nodos)
      \begin{scope}[on background layer]
        \node[grupoDiag, inner xsep=20pt, inner ysep=16pt, fit=(web)(grados)(chunks)(f0medida), label={[tituloGrupo]north:Fase 0 --- Colección de datos}] (f0) {};
        \node[grupoDiag, inner xsep=20pt, inner ysep=16pt, fit=(emb)(bd)(f1medida), label={[tituloGrupo]north:Fase 1 --- Indexación}] (f1) {};
        \node[grupoDiag, inner xsep=20pt, inner ysep=16pt, fit=(preg)(rec)(llm)(resp)(f2medida), label={[tituloGrupo]north:Fase 2 --- \textit{Pipeline} RAG}] (f2) {};
        \node[grupoDiag, inner xsep=20pt, inner ysep=16pt, fit=(chat)(f3medida), label={[tituloGrupo]north:Fase 3 --- Aplicación web}] (f3) {};
      \end{scope}
      % Flechas internas
      \draw[flechaDiag] (web) -- node[right, font=\footnotesize]{\textit{spider} (Scrapy)} (grados);
      \draw[flechaDiag] (grados) -- node[right, font=\footnotesize]{\textit{chunker}} (chunks);
      \draw[flechaDiag] (emb) -- (bd);
      \draw[flechaDiag] (preg) -- (rec);
      \draw[flechaDiag] (rec) -- (llm);
      \draw[flechaDiag] (llm) -- (resp);
      % Flechas entre fases
      \draw[flechaDiag] (chunks.east) -- (emb.west);
      \draw[flechaDiag] (bd.east) -- (rec.west);
      \draw[flechaDiag] (resp.east) -- (chat.south west);
    \end{tikzpicture}%
  }
  \caption{Arquitectura general del sistema, organizada por fases de
    desarrollo. Fuente: Elaboración propia.}
  \label{fig:arquitectura_pipeline}
\end{figure}
